
% Table to summarize types, traits, and operations

\subsection{Traits}
% Description, Purpose, and Motivation: What is this construct for? Why do we need it? What is it based off of (find source)?
% Semantics: What do traits represent? What does it mean when a type has a trait?
% For each trait:
% - Description, applications
% - Operations affected/unlocked by this trait
% - Assumptions, Invariants/Laws
% Maybe a table to summarize each trait?

\subsection{Boolean}
\subsection{Integer}
\subsection{Float}
\subsection{String}
\subsection{Set}
\subsection{Relation}
\subsection{Sequence}
\subsection{Bag}
\subsection{Tuple}
\subsection{Record}
\subsection{Procedure}
\subsection{Generic}
% Description, Purpose, and Motivation: What is this type for? Why do we need it? What is it based off of (find source)?
% Semantics: What do values represent? What does it mean to have an object of this type?
% Making this type: Example of its construction/constructors
% Examples of its use
% Relationships to other types (super/subtypes)
% Allowed traits (and how they affect this type)
% Restricted traits (and why we restrict them; what are the alternatives?)
% Formal definition
% Assumptions, Invariants/Laws
% Any other notes (edge cases, differences from normal execution)
\subsubsection{Operations}
% For each operation (include basic operations like equality, etc.):
% - Description: What is it used for? What is its purpose? Why do we need it? Where does it come from (find source)?
% - Example usage
% - Traits that affect this operator (and how they affect the operator)
% - Formalized type definition (with and without traits)
% - Any other notes/Detailed knowledge on this operator (edge cases, differences from normal execution)
% - Assumptions, Invariants/Laws